\chapter{Mecanismos de reacción de los complejos}\label{ch:b3:complex-mechanisms}

Si se disuelve una sal de cromo(III) en agua marcada con oxígeno-18, pasan días
antes de que el agua marcada haya sustituido a las moléculas de agua unidas al
metal; con una sal de cobre(II), el intercambio es completo en el tiempo de la
mezcla. La reacción es la misma, una molécula de agua que abandona un ion metálico y
otra que ocupa su lugar, y sin embargo su velocidad varía enormemente de un ion al
siguiente, y el número de electrones $d$ predice buena parte de esa variación. La manera
en que se sustituyen los ligandos decide cómo se fabrican los fármacos anticancerosos de platino y cómo
actúan; la manera en que los electrones pasan de un ion metálico a otro decide la velocidad de
la respiración y de toda batería. Este capítulo trata las dos familias de
reacciones de los complejos: la sustitución y la transferencia electrónica, esta última con
la teoría de Rudolph Marcus.

\begin{recall}
El volumen del segundo año describió los complejos, la denticidad de los ligandos, los ligandos
puente, los isómeros $fac$ y $mer$, el intercambio de ligandos como etapa elemental y
las energías de estabilización del campo cristalino con las configuraciones de alto y bajo espín;
el volumen del primer año, los mecanismos con las aproximaciones del estado estacionario y de la etapa
determinante de la velocidad. El \cref{ch:b3:rate-theories} dio la \omterm{thm:b3:rate-theories:eyring}{ecuación de Eyring} y el
significado de $\Delta^{\ddagger}S^\circ$; el \cref{ch:b3:electrode-kinetics}, la transferencia electrónica en un
electrodo, y el \cref{ch:b3:complex-spectra-magnetism}, los \omterm{def:b3:complex-spectra-magnetism:ligand-field-term}{términos del campo de ligandos} y el
efecto Jahn--Teller.
\end{recall}

\section{Labilidad e inercia}

\begin{definition}[Complejos lábiles e inertes]\label{def:b3:complex-mechanisms:labile}
Un complejo es \emph{lábil}\index{complejo labil@complejo lábil} si sus ligandos se sustituyen
deprisa, y un \emph{complejo cinéticamente inerte}\index{complejo cineticamente inerte@complejo cinéticamente inerte}
si se sustituyen lentamente; siguiendo a Taube, se dice que un complejo es lábil cuando
sus reacciones con un ligando a \qty{0.1}{mol/L} y \qty{25}{\celsius} terminan en un
minuto aproximadamente. La labilidad es una propiedad cinética, sin relación con la estabilidad
termodinámica del complejo.
\end{definition}

El intercambio de agua, $\mathrm{[M(H_2O)_6]^{n+} + H_2O^{\ast} \to [M(H_2O)_5(H_2O^{\ast})]^{n+} + H_2O}$, es la sustitución
más sencilla y fija una escala. Los iones de los grupos principales intercambian más deprisa cuanto
mayores y menos cargados son. Entre los iones de los metales de transición, los de configuración $d^3$
(\ce{Cr^{3+}}) o $d^6$ de bajo espín (\ce{Co^{3+}}, \ce{Rh^{3+}}, \ce{Ir^{3+}}) son inertes, como
la mayoría de los metales de la segunda y la tercera fila, mientras que $d^9$ (\ce{Cu^{2+}}) y $d^4$ de alto espín (\ce{Cr^{2+}}),
distorsionados por el efecto Jahn--Teller, son extremadamente \omterm{def:b3:complex-mechanisms:labile}{lábiles}.

\begin{proposition}[Energía de activación del campo de ligandos]\label{prop:b3:complex-mechanisms:lfae}
Si una sustitución pasa por un estado de transición pentacoordinado de pirámide de base
cuadrada, la pérdida de estabilización del campo de ligandos al llegar a él (la energía de
activación del campo de ligandos) es, en el modelo más sencillo, solo $\sigma$, máxima para $d^6$ de bajo espín
($4\,Dq$), y después para $d^3$ y $d^8$ ($2\,Dq$), nula para $d^0$, $d^5$ de alto espín y $d^{10}$, y negativa
para $d^4$ y $d^9$ de alto espín.
\end{proposition}

\begin{proof}
Se modeliza cada ligando como si elevara un orbital $d$ en $e_\sigma$ veces el cuadrado de la
amplitud angular del orbital en su dirección, normalizada a 1 sobre el eje del lóbulo: un ligando
axial da a $d_{z^2}$ 1, uno ecuatorial da a $d_{z^2}$ $\frac14$ y a $d_{x^2-y^2}$ $\frac34$; los orbitales $t_{2g}$ apuntan
entre los ligandos y no reciben nada. El octaedro da a $d_{z^2}$ y a $d_{x^2-y^2}$ $3e_\sigma$ cada uno
($\Delta_{\mathrm o} = 3e_\sigma = 10Dq$); la pirámide de base cuadrada, sin uno de los ligandos axiales, da a $d_{z^2}$ $2e_\sigma$ y a
$d_{x^2-y^2}$ $3e_\sigma$. La estabilización de $n$ electrones es su energía menos $n/5$ veces la suma de
los cinco niveles (el baricentro). Para $d^6$ de bajo espín ($t_{2g}^6$): $0 - \frac65 \times 6e_\sigma$ en el
octaedro, $0 - \frac65 \times 5e_\sigma$ en la pirámide: una pérdida de $1.2e_\sigma = 4Dq$. Para $d^3$: $\frac35e_\sigma =
2Dq$; para $d^8$ ($t_{2g}^6e_g^2$): $(5 - \frac85 \times 5) - (6 - \frac85 \times 6) = 0.6e_\sigma$. Para $d^9$: $(7 - 9) - (9 -
10.8) = -0.2e_\sigma$.
\end{proof}

\begin{omfigure}
\begin{tikzpicture}
\begin{axis}[width=0.72\linewidth, height=5.4cm, xmin=-0.6, xmax=10.6, ymin=-1.2, ymax=4.6,
  axis x line=bottom, axis y line=left, ycomb, xtick={0,1,2,3,4,5,6,7,8,9,10},
  xlabel={$n$ ($d^n$)}, ylabel={energía de activación / $Dq$}, tick label style={font=\small},
  label style={font=\small}]
\draw[black!40] (axis cs:-0.6,0) -- (axis cs:10.6,0);
\addplot[omDef, line width=4pt, xshift=-2.5pt] table[x=n, y=hs] {figdata/out/bachelor-3-ligand-field-activation.dat};
\addplot[omThm, line width=4pt, xshift=2.5pt, unbounded coords=discard] table[x=n, y=ls] {figdata/out/bachelor-3-ligand-field-activation.dat};
\end{axis}
\end{tikzpicture}

{\small Energías de activación del campo de ligandos para un camino disociativo a través de una
pirámide de base cuadrada, en el modelo solo $\sigma$ de la proposición (azul: alto espín;
rojo: bajo espín). Los valores mayores,
para $d^6$ de bajo espín, $d^3$ y $d^8$, corresponden a los iones inertes o lentos (\ce{Co^{3+}}, \ce{Cr^{3+}}, \ce{Ni^{2+}});
los negativos, para $d^4$ y $d^9$, a los iones de Jahn--Teller muy \omterm{def:b3:complex-mechanisms:labile}{lábiles}.}
\end{omfigure}

\begin{definition}[Volumen de activación]\label{def:b3:complex-mechanisms:activation-volume}
El \emph{volumen de activación}\index{volumen de activacion@volumen de activación} $\Delta V^{\ddagger}$ de una etapa
elemental es la diferencia entre el volumen molar parcial del \omterm{def:b3:rate-theories:activated-complex}{complejo
activado} y los de los reactivos.
\end{definition}

\begin{proposition}[Dependencia de una constante de velocidad con la presión]\label{prop:b3:complex-mechanisms:pressure}
A temperatura constante, $\displaystyle\Bigl(\frac{\partial\ln k}{\partial p}\Bigr)_T = -\frac{\Delta V^{\ddagger}}{RT}$.
\end{proposition}

\begin{proof}
Por la \omterm{thm:b3:rate-theories:eyring}{ecuación de Eyring}, $\ln k = \text{cte} - \Delta^{\ddagger}G/RT$, y $(\partial\Delta^{\ddagger}G/\partial p)_T = \Delta V^{\ddagger}$, como para
cualquier energía de Gibbs.
\end{proof}

Una etapa en la que un ligando se va tiene $\Delta V^{\ddagger} > 0$; una en la que un ligando entra tiene $\Delta V^{\ddagger}
< 0$. Medido a unos cientos de megapascales, el \omterm{def:b3:complex-mechanisms:activation-volume}{volumen de activación} es el diagnóstico más
directo de un mecanismo de sustitución, menos enturbiado por la solvatación que
la \omterm{def:b3:rate-theories:activation}{entropía de activación}.

\section{Mecanismos de sustitución}

\begin{definition}[Mecanismos disociativo, asociativo y de intercambio]\label{def:b3:complex-mechanisms:mechanisms}
En un \emph{mecanismo disociativo}\index{mecanismo disociativo} (D), el ligando
saliente se va primero, lo que da un intermedio de número de coordinación menor; en un
\emph{mecanismo asociativo}\index{mecanismo asociativo} (A), el ligando entrante
se une primero, lo que da un intermedio de número de coordinación mayor; en un
\emph{mecanismo de intercambio}\index{mecanismo de intercambio} (I), los dos ligandos
se intercambian en una sola etapa, con carácter de ruptura de enlace ($I_d$) o de formación de enlace ($I_a$).
\end{definition}

\begin{proposition}[Leyes de velocidad de los mecanismos D y A]\label{prop:b3:complex-mechanisms:d-a-laws}
Para $\mathrm{ML_5X + Y \to ML_5Y + X}$: un mecanismo D ($\mathrm{ML_5X} \rightleftharpoons \mathrm{ML_5} + \mathrm X$, $k_1$, $k_{-1}$; $\mathrm{ML_5} + \mathrm Y
\to \mathrm{ML_5Y}$, $k_2$) da
\[
  v = \frac{k_1k_2[\mathrm{ML_5X}][\mathrm Y]}{k_{-1}[\mathrm X] + k_2[\mathrm Y]},
\]
de primer orden respecto del complejo e independiente de $[\mathrm Y]$ a $[\mathrm Y]$ alta; un mecanismo A
a través de un intermedio heptacoordinado en estado estacionario da $v =
k[\mathrm{ML_5X}][\mathrm Y]$, de primer orden respecto de cada uno.
\end{proposition}

\begin{proof}
D: estado estacionario para $\mathrm{ML_5}$, $k_1[\mathrm{ML_5X}] = (k_{-1}[\mathrm X] + k_2[\mathrm Y])[\mathrm{ML_5}]$, y $v = k_2[\mathrm{ML_5}][\mathrm Y]$.
A $[\mathrm Y]$ alta, $v \to k_1[\mathrm{ML_5X}]$. A: $\mathrm{ML_5X} + \mathrm Y \rightleftharpoons \mathrm{ML_5XY}$ ($k_a$, $k_{-a}$), $\mathrm{ML_5XY} \to
\mathrm{ML_5Y} + \mathrm X$ ($k_b$); el estado estacionario da $v = \frac{k_ak_b}{k_{-a} + k_b}[\mathrm{ML_5X}][\mathrm Y]$.
\end{proof}

En agua, el ligando entrante suele estar presente a una concentración mucho menor
que el disolvente, y la verdadera primera etapa es el encuentro del complejo con el
ligando.

\begin{definition}[Mecanismo de Eigen--Wilkins]\label{def:b3:complex-mechanisms:eigen-wilkins}
El \emph{mecanismo de Eigen--Wilkins}\index{mecanismo de Eigen--Wilkins} de
sustitución de un acuocomplejo es un preequilibrio rápido que forma un
\emph{complejo de esfera externa}\index{complejo de esfera externa}, en el que el ligando
entrante se sitúa junto al complejo, fuera de su esfera de coordinación, seguido del
intercambio, determinante de la velocidad, de una molécula de agua por ese ligando.
\end{definition}

\begin{theorem}[Ley de velocidad de Eigen--Wilkins]\label{thm:b3:complex-mechanisms:eigen-wilkins}
Con la constante de equilibrio de esfera externa $K_{\mathrm{os}}$ y la constante de velocidad de intercambio
$k_i$, la constante de velocidad de pseudoprimer orden observada con $[\mathrm Y]$ en exceso es
\[
  k_{\mathrm{obs}} = \frac{K_{\mathrm{os}}k_i[\mathrm Y]}{1 + K_{\mathrm{os}}[\mathrm Y]} .
\]
\end{theorem}

\begin{proof}
El complejo se reparte entre la forma libre y la forma de esfera externa, $[\mathrm{M\cdot Y}] = K_{\mathrm{os}}[\mathrm M][\mathrm Y]$,
con un total $[\mathrm M]_{\mathrm{tot}} = [\mathrm M](1 + K_{\mathrm{os}}[\mathrm Y])$; la velocidad es $k_i[\mathrm{M\cdot Y}] = k_iK_{\mathrm{os}}[\mathrm Y][\mathrm M]_{\mathrm{tot}}/(1 + K_{\mathrm{os}}[\mathrm Y])$.
\end{proof}

A $[\mathrm Y]$ baja, la reacción es de segundo orden con $k = K_{\mathrm{os}}k_i$; $K_{\mathrm{os}}$, estimada a partir de las
cargas y de la distancia de máxima aproximación, es casi la misma para todos los
ligandos de la misma carga, de modo que $k_i$, la etapa de intercambio de agua, controla la
velocidad: la sustitución sobre un acuoion es casi independiente del ligando
entrante, la firma de un intercambio disociativo.

\begin{definition}[Mecanismo de la base conjugada]\label{def:b3:complex-mechanisms:conjugate-base}
El \emph{mecanismo de la base conjugada}\index{mecanismo de la base conjugada} de la hidrólisis
básica de un complejo ammínico, $\mathrm{[Co(NH_3)_5X]^{2+} + OH^- \to [Co(NH_3)_5(OH)]^{2+} + X^-}$, es
la desprotonación de un ligando ammina en un preequilibrio rápido, seguida de la
pérdida disociativa de $\mathrm X^-$ desde la base conjugada amido, cuyo ligando $\mathrm{NH_2^-}$
labiliza el complejo.
\end{definition}

\begin{proposition}[Ley de velocidad de la hidrólisis básica]\label{prop:b3:complex-mechanisms:conjugate-base}
Con la constante de desprotonación $K$ y la constante de velocidad de disociación $k$ de la
base conjugada, $v = \dfrac{kK[\text{complejo}]_{\mathrm{tot}}[\ce{OH-}]}{1 + K[\ce{OH-}]}$, de segundo orden a $[\ce{OH-}]$ baja.
\end{proposition}

\begin{proof}
Como para la ley de Eigen--Wilkins: el complejo se reparte entre su forma ácida y su
forma de base conjugada en la razón $1 : K[\ce{OH-}]$, y solo la segunda reacciona.
\end{proof}

La ley de velocidad sola no distingue este mecanismo de un ataque directo del
\ce{OH-}; la prueba es que la hidrólisis básica necesita un protón N--H (los complejos que no lo
tienen no la presentan) y que los protones de los ligandos ammina se intercambian con el disolvente mucho
más deprisa que la hidrólisis.

\begin{method}[Diagnosticar un mecanismo de sustitución]\label{met:b3:complex-mechanisms:diagnose}
\begin{enumerate}
  \item Dependencia del grupo entrante: velocidad independiente de Y (o con saturación),
        disociativo; velocidad proporcional a $[\mathrm Y]$ y sensible a la naturaleza de Y,
        asociativo.
  \item \omterm{def:b3:complex-mechanisms:activation-volume}{Volumen de activación}: positivo, disociativo; negativo, asociativo.
  \item \omterm{def:b3:rate-theories:activation}{Entropía de activación}: positiva para D, negativa para A (con cautela, ya que
        la solvatación contribuye).
  \item Dependencia del grupo saliente: una velocidad que sigue la fuerza del enlace M--X
        apunta a una ruptura de enlace en el estado de transición.
\end{enumerate}
\end{method}

\begin{omfigure}
\begin{tikzpicture}
\begin{axis}[name=pd, omprofile, width=0.46\linewidth, height=4.6cm, xmin=0, xmax=10, ymin=0, ymax=10,
  xlabel={coordenada de reacción}, ylabel={energía}, title={\small D: $\mathrm{ML_5X} \to \mathrm{ML_5} + \mathrm X \to \mathrm{ML_5Y}$}]
\addplot[omDef, very thick, smooth] coordinates {(0,2) (1.5,2.1) (3,7.5) (4.2,5.6) (5.0,5.4) (5.8,5.6) (7,7.0) (8.5,1.6) (10,1.5)};
\node[font=\scriptsize, anchor=north, align=center] at (axis cs:5,5.2) {ML$_5$ (NC 5)};
\end{axis}
\begin{axis}[at={(pd.east)}, anchor=west, xshift=1.0cm, omprofile, width=0.46\linewidth, height=4.6cm,
  xmin=0, xmax=10, ymin=0, ymax=10, xlabel={coordenada de reacción}, ylabel={energía},
  title={\small A: $\mathrm{ML_5X} + \mathrm Y \to \mathrm{ML_5XY} \to \mathrm{ML_5Y}$}]
\addplot[omThm, very thick, smooth] coordinates {(0,2) (1.5,2.1) (3,6.5) (4.2,4.4) (5.0,4.2) (5.8,4.4) (7,7.2) (8.5,1.6) (10,1.5)};
\node[font=\scriptsize, anchor=north, align=center] at (axis cs:5,4.0) {ML$_5$XY (NC 7)};
\end{axis}
\end{tikzpicture}

{\small Perfiles de energía de una sustitución disociativa y de una asociativa
(esquemáticos): ambas pasan por un intermedio, de número de coordinación menor
(NC) en D y de número de coordinación mayor en A. En los \omterm{def:b3:complex-mechanisms:mechanisms}{mecanismos de intercambio},
el pozo desaparece y solo queda un estado de transición.}
\end{omfigure}

Un complejo octaédrico que se sustituye a través de una pirámide de base cuadrada conserva
en general su configuración ($cis$ sigue siendo $cis$); uno que pasa por una bipirámide
trigonal, como hacen muchos complejos de cobalto(III) en la hidrólisis básica, puede dar una
mezcla de productos $cis$ y $trans$.

\section{Sustitución plano-cuadrada y efecto trans}

Los complejos plano-cuadrados de iones $d^8$ (\ce{Pt^{2+}}, \ce{Pd^{2+}}, \ce{Au^{3+}}, \ce{Ni^{2+}} con ligandos fuertes)
están coordinativamente insaturados: un ligando entrante puede aproximarse por la
dirección axial vacía. Sus sustituciones son asociativas, a través de un
estado de transición pentacoordinado de bipirámide trigonal.

\begin{proposition}[Ley de velocidad de dos términos]\label{prop:b3:complex-mechanisms:square-planar}
La sustitución $\mathrm{[PtL_3X] + Y \to [PtL_3Y] + X}$ en un disolvente coordinante S sigue
$k_{\mathrm{obs}} = k_1 + k_2[\mathrm Y]$ con Y en exceso.
\end{proposition}

\begin{proof}
Dos caminos asociativos paralelos: el ataque directo de Y ($k_2[\mathrm Y]$) y el ataque del
disolvente, presente en gran exceso (pseudoprimer orden, $k_1$), que da $[\mathrm{PtL_3S}]$, que
Y transforma después rápidamente. Los caminos paralelos de primer orden se suman.
\end{proof}

\begin{definition}[Efecto trans, influencia trans]\label{def:b3:complex-mechanisms:trans-effect}
El \emph{efecto trans}\index{efecto trans} es la aceleración de la sustitución
de un ligando por el ligando situado en trans respecto de él: un efecto cinético, sobre el estado de transición.
La \emph{influencia trans}\index{influencia trans} es el debilitamiento, en el estado
fundamental, del enlace en trans respecto de un ligando, visible en las longitudes de enlace y en las frecuencias
de vibración: un efecto termodinámico.
\end{definition}

El \omterm{def:b3:complex-mechanisms:trans-effect}{efecto trans} sigue el orden $\ce{CO}$, $\ce{CN-}$, $\ce{C2H4} > \ce{PR3}$, $\ce{H-} > \ce{CH3-} > \ce{I-}$, $\ce{SCN-} >
\ce{Br-} > \ce{Cl-} > \ce{NH3}$, $\ce{H2O}$. Los $\sigma$-dadores fuertes debilitan el enlace en trans (\omterm{def:b3:complex-mechanisms:trans-effect}{influencia trans});
los $\pi$-aceptores fuertes estabilizan el estado de transición pentacoordinado al retirar
\omterm{def:b3:computational-chemistry:dft}{densidad electrónica} del metal. Ambos aumentan la velocidad.

\begin{method}[Planificar la síntesis de isómeros de platino(II)]\label{met:b3:complex-mechanisms:pt-synthesis}
\begin{enumerate}
  \item En cada etapa, el ligando sustituido es el que está en trans respecto del ligando de
        mayor \omterm{def:b3:complex-mechanisms:trans-effect}{efecto trans} presente.
  \item Ordenar las adiciones de modo que esta regla conduzca al isómero deseado.
\end{enumerate}
\end{method}

\begin{omfigure}
\setchemfig{atom sep=2.2em}
\schemestart
\chemfig{Cl-Pt(-[2]Cl)(-[6]Cl)-Cl}
\arrow{->[\ce{NH3}]}
\chemfig{Cl-Pt(-[2]Cl)(-[6]Cl)-NH_3}
\arrow{->[\ce{NH3}]}
\chemfig{Cl-Pt(-[2]NH_3)(-[6]Cl)-NH_3}
\schemestop

\bigskip
\schemestart
\chemfig{H_3N-Pt(-[2]NH_3)(-[6]NH_3)-NH_3}
\arrow{->[\ce{Cl-}]}
\chemfig{H_3N-Pt(-[2]NH_3)(-[6]NH_3)-Cl}
\arrow{->[\ce{Cl-}]}
\chemfig{Cl-Pt(-[2]NH_3)(-[6]NH_3)-Cl}
\schemestop

{\small El \omterm{def:b3:complex-mechanisms:trans-effect}{efecto trans} en acción (se omiten las cargas). Arriba: a partir de $\ce{[PtCl4]^{2-}}$, el segundo
amoníaco sustituye a un cloruro situado en trans respecto de un cloruro (Cl$^-$ tiene el mayor \omterm{def:b3:complex-mechanisms:trans-effect}{efecto trans}),
lo que da el isómero $cis$, el cisplatino. Abajo: a partir de $\ce{[Pt(NH3)4]^{2+}}$, el segundo cloruro
sustituye al amoníaco situado en trans respecto del primer cloruro, lo que da el isómero $trans$.}
\end{omfigure}

\section{Transferencia electrónica}

\begin{definition}[Transferencia electrónica de esfera externa y de esfera interna]\label{def:b3:complex-mechanisms:electron-transfer}
En la \emph{transferencia electrónica de esfera externa}\index{transferencia electronica de esfera externa@transferencia electrónica de esfera externa}, el
electrón pasa entre dos complejos cuyas esferas de coordinación permanecen intactas;
en la \emph{transferencia electrónica de esfera interna}\index{transferencia electronica de esfera interna@transferencia electrónica de esfera interna},
pasa a través de un ligando que tiende un puente entre los dos metales en un complejo binuclear
transitorio. Una \emph{reacción de autointercambio}\index{reaccion de autointercambio@reacción de autointercambio} es una
transferencia electrónica entre los dos estados de oxidación de un mismo par, sin
cambio químico neto, como $\ce{[Fe(H2O)6]^{3+} + [Fe(H2O)6]^{2+}}$.
\end{definition}

Henry Taube demostró el camino de esfera interna en 1953 con una reacción elegida de modo que
el producto contara su propia historia. El cobalto(III) es inerte, y el cobalto(II), \omterm{def:b3:complex-mechanisms:labile}{lábil}; el cromo(II)
es \omterm{def:b3:complex-mechanisms:labile}{lábil}, y el cromo(III), inerte. Cuando el $\ce{[Co(NH3)5Cl]^{2+}}$ se reduce con $\ce{[Cr(H2O)6]^{2+}}$ en
medio ácido, todo el cromo(III) formado lleva el cloruro, como $\ce{[Cr(H2O)5Cl]^{2+}}$, incluso en
presencia de cloruro libre: el cloruro debió de tender un puente entre los dos metales,
quedar unido al cromo(III) inerte en cuanto pasó el electrón y ser
liberado por el cobalto(II), \omterm{def:b3:complex-mechanisms:labile}{lábil}.

\begin{omfigure}
\begin{tikzpicture}[font=\small]
  \begin{scope}[scale=0.8]
    \pic[transform shape] (a) at (0,0) {omoctahedron};
    \pic[transform shape] (b) at (2.0,0) {omoctahedron};
  \end{scope}
  \node[anchor=north east] at (0.3,-1.0) {$\mathrm{Co^{III}(NH_3)_5}$};
  \node[anchor=north west] at (1.3,-1.0) {$\mathrm{Cr^{II}(H_2O)_5}$};
  \node[omThm, anchor=south] at (0.8,0.15) {Cl};
  \draw[->, thick, omDef] (1.6,1.75) to[bend right=40] node[midway, above] {e$^-$} (0,1.75);
\end{tikzpicture}

{\small El intermedio con puente de Taube (esquema): el cloruro, todavía unido al
cobalto(III), entra en la esfera de coordinación del cromo(II) sustituyendo a una molécula de
agua; el electrón pasa a través del puente, y el cloruro se queda en el
cromo(III), ahora inerte.}
\end{omfigure}

Un electrón salta en unos $10^{-16}$ s, mucho más deprisa de lo que se mueven los núcleos: según el
\omterm{thm:b3:electronic-spectroscopy:franck-condon}{principio de Franck--Condon} (\cref{ch:b3:electronic-spectroscopy}), la transferencia
solo puede producirse en una configuración nuclear en la que reactivos y productos tienen
la misma energía. En el autointercambio del hierro, los enlaces \ce{Fe^{III}}--O son más cortos que
los enlaces \ce{Fe^{II}}--O; antes de que el electrón pueda moverse, los dos complejos deben distorsionarse hasta
una geometría intermedia común, con un coste energético.

\section{La teoría de Marcus}

\begin{definition}[Energía de reorganización]\label{def:b3:complex-mechanisms:reorganisation}
La \emph{energía de reorganización}\index{energia de reorganizacion@energía de reorganización} $\lambda$ de una
transferencia electrónica es la energía necesaria para llevar los reactivos, sin
transferir el electrón, a la configuración nuclear (longitudes de enlace de los
complejos y orientaciones del disolvente circundante) de los productos. Es la
suma de una parte interna, debida a los enlaces, y de una parte externa, debida al disolvente.
\end{definition}

\begin{theorem}[Ecuación de Marcus]\label{thm:b3:complex-mechanisms:marcus}
Si las energías de Gibbs de reactivos y productos son parábolas de la misma
curvatura en una coordenada de reacción $x$ (0 en el equilibrio de los reactivos, 1 en el de los
productos), la \omterm{def:b3:rate-theories:activation}{energía de Gibbs de activación} de una transferencia electrónica con
energía de Gibbs estándar de reacción $\Delta G^\circ$ y \omterm{def:b3:complex-mechanisms:reorganisation}{energía de reorganización} $\lambda$ es
\[
  \Delta G^{\ddagger} = \frac{(\lambda + \Delta G^\circ)^2}{4\lambda} .
\]
Es la \emph{ecuación de Marcus}\index{ecuacion de Marcus@ecuación de Marcus}.
\end{theorem}

\begin{proof}
Escríbase $G_R = \lambda x^2$ y $G_P = \lambda(x - 1)^2 + \Delta G^\circ$: la curvatura queda fijada por $G_R(1) = \lambda$, la
definición de $\lambda$. La transferencia se produce donde las curvas se cortan: $\lambda x^2 = \lambda x^2 - 2\lambda x + \lambda +
\Delta G^\circ$, luego $x^\ast = (\lambda + \Delta G^\circ)/2\lambda$, y $\Delta G^{\ddagger} = G_R(x^\ast) = (\lambda + \Delta G^\circ)^2/4\lambda$.
\end{proof}

\begin{corollary}[Región invertida]\label{cor:b3:complex-mechanisms:inverted}
Para $\lambda$ fija, la velocidad aumenta a medida que la reacción se hace más exergónica hasta
$-\Delta G^\circ = \lambda$, donde $\Delta G^{\ddagger} = 0$; más allá, vuelve a disminuir.
\end{corollary}

\begin{proof}
$\Delta G^{\ddagger}$ es una parábola en $\Delta G^\circ$ con su mínimo, cero, en $\Delta G^\circ = -\lambda$; para $\Delta G^\circ < -\lambda$ vuelve
a crecer.
\end{proof}

\begin{definition}[Región invertida]\label{def:b3:complex-mechanisms:inverted}
La \emph{región invertida}\index{region invertida@región invertida} de la transferencia electrónica es el
intervalo de fuerzas impulsoras $-\Delta G^\circ > \lambda$ en el que la velocidad disminuye a medida que la reacción se hace
más exergónica.
\end{definition}

\begin{omfigure}
\begin{tikzpicture}
\begin{axis}[name=mp, width=0.47\linewidth, height=5.6cm, xmin=-0.6, xmax=2.2, ymin=-160, ymax=120,
  axis lines=left, xlabel={coordenada de reacción $x$}, ylabel={$G$ / \unit{kJ.mol^{-1}}},
  tick label style={font=\small}, label style={font=\small}, xtick={0,1,2}]
\addplot[black, thick] table[x=x, y=GR] {figdata/out/bachelor-3-marcus-parabolas.dat};
\addplot[omDef!50, thick] table[x=x, y=P0] {figdata/out/bachelor-3-marcus-parabolas.dat};
\addplot[omDef, thick] table[x=x, y=P50] {figdata/out/bachelor-3-marcus-parabolas.dat};
\addplot[omThm, thick] table[x=x, y=P100] {figdata/out/bachelor-3-marcus-parabolas.dat};
\addplot[omThm!60, thick, dashed] table[x=x, y=P150] {figdata/out/bachelor-3-marcus-parabolas.dat};
\node[font=\scriptsize, anchor=north] at (axis cs:0,-2) {R};
\end{axis}
\begin{axis}[at={(mp.east)}, anchor=west, xshift=1.4cm, width=0.47\linewidth, height=5.6cm,
  xmin=0, xmax=250, ymin=0, ymax=11.5, axis lines=left, xlabel={$-\Delta G^\circ$ / \unit{kJ.mol^{-1}}},
  ylabel={$\log(k/\unit{s^{-1}})$}, tick label style={font=\small}, label style={font=\small}]
\addplot[omDef, very thick] table[x=mdG, y=logk] {figdata/out/bachelor-3-marcus-rate.dat};
\draw[black!40, dashed] (axis cs:100,0) -- (axis cs:100,11);
\node[font=\scriptsize, anchor=south east] at (axis cs:95,1) {normal};
\node[font=\scriptsize, anchor=south west] at (axis cs:105,1) {invertida};
\end{axis}
\end{tikzpicture}

{\small Teoría de Marcus con $\lambda = \qty{100}{kJ/mol}$ (modelo). Izquierda: la parábola de los reactivos
(negro) y las parábolas de los productos para $\Delta G^\circ = 0$ y $-50$ (azul), $-100 = -\lambda$ (rojo) y
\qty{-150}{kJ/mol} (trazo discontinuo, \omterm{def:b3:complex-mechanisms:inverted}{región invertida}); su punto de
corte, el estado de transición, baja hasta el fondo de la curva de los reactivos en
$-\Delta G^\circ = \lambda$ y vuelve a subir más allá. Derecha: el logaritmo de la constante de velocidad
(factor preexponencial del modelo \qty{e11}{s^{-1}}) pasa por un máximo en $-\Delta G^\circ = \lambda$.}
\end{omfigure}

\begin{proposition}[Energía de reorganización de esfera externa]\label{prop:b3:complex-mechanisms:outer-reorganisation}
Para dos reactivos esféricos de radios $a_1$, $a_2$ a una distancia $d$ en un disolvente de
índice de refracción $n$ y permitividad relativa $\varepsilon_s$, la parte de $\lambda$ debida al disolvente es
proporcional a $\bigl(\frac{1}{2a_1} + \frac{1}{2a_2} - \frac1d\bigr)\bigl(\frac{1}{n^2} - \frac{1}{\varepsilon_s}\bigr)$.
\end{proposition}

\admitted

El agua, muy polar ($\varepsilon_s$ grande) pero con un índice de refracción corriente, da una
reorganización del disolvente grande; los disolventes apolares, una pequeña. Los reactivos grandes
se reorganizan menos: por eso las proteínas de transferencia electrónica entierran sus sitios metálicos
en el interior de la proteína.

\begin{theorem}[Relación cruzada de Marcus]\label{thm:b3:complex-mechanisms:cross-relation}
Para una reacción cruzada $\mathrm{A_{ox} + B_{red} \to A_{red} + B_{ox}}$ de constante de equilibrio $K_{12}$, entre pares
cuyas constantes de velocidad de autointercambio son $k_{11}$ y $k_{22}$,
\[
  k_{12} = \sqrt{k_{11}k_{22}K_{12}f}, \qquad \ln f = \frac{(\ln K_{12})^2}{4\ln(k_{11}k_{22}/Z^2)},
\]
donde $Z$ es el factor de frecuencia de colisión; $f \approx 1$ cuando $K_{12}$ no es demasiado grande. Es la
\emph{relación cruzada de Marcus}\index{relacion cruzada de Marcus@relación cruzada de Marcus}.
\end{theorem}

\begin{proof}[Demostración parcial]
Supóngase que la \omterm{def:b3:complex-mechanisms:reorganisation}{energía de reorganización} de la reacción cruzada es la media de las de los
autointercambios, $\lambda_{12} = (\lambda_{11} + \lambda_{22})/2$ (cada reactivo aporta la mitad de la reorganización de su
propio autointercambio), y que todas las constantes de velocidad son $k = Z\eu^{-\Delta G^{\ddagger}/RT}$. Entonces
$\Delta G_{12}^{\ddagger} = \frac{\lambda_{12}}{4}\bigl(1 + \frac{\Delta G^\circ_{12}}{\lambda_{12}}\bigr)^2 = \frac{\lambda_{12}}{4} + \frac{\Delta G^\circ_{12}}{2} + \frac{(\Delta G^\circ_{12})^2}{4\lambda_{12}}$. El primer término da
$\sqrt{k_{11}k_{22}}$ (ya que $\Delta G^{\ddagger}_{ii} = \lambda_{ii}/4$), el segundo, $\sqrt{K_{12}}$ (con $\Delta G^\circ_{12} = -RT\ln K_{12}$), y el tercero,
el factor $f$.
\end{proof}

\begin{method}[La velocidad de una reacción cruzada a partir de datos de autointercambio]\label{met:b3:complex-mechanisms:marcus-estimate}
\begin{enumerate}
  \item Encontrar las constantes de autointercambio $k_{11}$, $k_{22}$ de los dos pares.
  \item Calcular $K_{12}$ a partir de los potenciales estándar: $\ln K_{12} = nF(E^\circ_{\text{oxidante}} - E^\circ_{\text{reductor}})/RT$.
  \item $k_{12} = \sqrt{k_{11}k_{22}K_{12}}$ (con $f$ si $K_{12}$ es muy grande). Una concordancia con el valor medido
        dentro de un factor de pocas unidades indica un mecanismo de esfera externa.
\end{enumerate}
\end{method}

Marcus publicó la teoría en 1956; la \omterm{def:b3:complex-mechanisms:inverted}{región invertida}, puesta en duda al principio, fue
observada en 1984 por Gerhard Closs y John Miller en moléculas en las que el dador y el
aceptor se mantienen a una distancia fija mediante un espaciador rígido, de modo que la difusión
no puede ocultarla.

\begin{inthelab}[El intercambio de agua por RMN de oxígeno-17]
La resonancia del oxígeno-17 del agua unida a un ion paramagnético y la del agua
libre se ensanchan por el intercambio entre ambas. Registrando la anchura de raya de la
señal del agua libre en función de la temperatura, en una disolución del perclorato metálico
enriquecida en \ce{^{17}O}, se obtiene la constante de velocidad de intercambio y, con una sonda de alta presión,
su \omterm{def:b3:complex-mechanisms:activation-volume}{volumen de activación}.
\end{inthelab}

\begin{safety}{\ghs[0.8cm]{GHS05}\ghs[0.8cm]{GHS06}\ghs[0.8cm]{GHS08}}
% ledger: ghs4:K2PtCl4
Tetracloroplatinato(II) de potasio: tóxico en caso de ingestión, provoca lesiones oculares
graves, sensibilizante respiratorio y cutáneo. Los complejos de platino se manipulan con
guantes en una vitrina; el cisplatino y sus análogos son fármacos citotóxicos y se
manipulan como tales.
\end{safety}

\begin{history}[Taube y Marcus]
Henry Taube clasificó los complejos en \omterm{def:b3:complex-mechanisms:labile}{lábiles} e inertes en 1952, y con su
experimento del cromo y el cobalto de 1953 estableció el mecanismo de esfera interna;
recibió el Premio Nobel de Química de 1983. Rudolph Marcus dedujo a partir de 1956 la
teoría de la transferencia electrónica que lleva su nombre, y recibió el Premio Nobel
de Química de 1992.
\end{history}

\section{Ejercicios}

\begin{exercise}[$\star$]\label{exo:b3:complex-mechanisms:1}
Clasifíquense como \omterm{def:b3:complex-mechanisms:labile}{lábiles} o inertes, justificándolo: $\ce{[Cr(H2O)6]^{3+}}$, $\ce{[Co(NH3)6]^{3+}}$, $\ce{[Cr(H2O)6]^{2+}}$,
$\ce{[Cu(H2O)6]^{2+}}$, $\ce{[Zn(H2O)6]^{2+}}$.
\end{exercise}

\begin{exercise}[$\star$]\label{exo:b3:complex-mechanisms:2}
Demuéstrese que la ley de velocidad D se reduce a $v = k_1[\mathrm{ML_5X}]$ a $[\mathrm Y]$ alta y a $v = (k_1k_2/k_{-1})
[\mathrm{ML_5X}][\mathrm Y]/[\mathrm X]$ cuando $k_{-1}[\mathrm X] \gg k_2[\mathrm Y]$.
\end{exercise}

\begin{exercise}[$\star$]\label{exo:b3:complex-mechanisms:3}
Predíganse los signos de $\Delta V^{\ddagger}$ y de $\Delta^{\ddagger}S^\circ$ para una sustitución D y para una A.
\end{exercise}

\begin{exercise}[$\star$]\label{exo:b3:complex-mechanisms:4}
Dense los productos del $\ce{[PtCl4]^{2-}}$ con dos \ce{NH3}, y del $\ce{[Pt(NH3)4]^{2+}}$ con dos \ce{Cl-}.
\end{exercise}

\begin{exercise}[$\star\star$]\label{exo:b3:complex-mechanisms:5}
Para una sustitución sobre un acuoion, $K_{\mathrm{os}} = \qty{2.0}{L.mol^{-1}}$ y $k_i = \qty{3.0e4}{s^{-1}}$ (datos del
ejercicio). Calcúlese $k_{\mathrm{obs}}$ para $[\mathrm Y] = 0.10$ y \qty{1.0}{mol/L}, y su límite.
\end{exercise}

\begin{exercise}[$\star\star$]\label{exo:b3:complex-mechanisms:6}
Una constante de velocidad se duplica cuando la presión sube de 0.1 a \qty{100}{MPa} a \qty{298}{K}.
Calcúlese $\Delta V^{\ddagger}$ y extráigase una conclusión.
\end{exercise}

\begin{exercise}[$\star\star$]\label{exo:b3:complex-mechanisms:7}
En el $\ce{[PtCl3(PEt3)]-}$, el enlace Pt--Cl en trans respecto de \ce{PEt3} mide \qty{2.42}{\angstrom}, y los otros dos, \qty{2.32}{\angstrom}
(datos del ejercicio). Interprétese y predígase qué cloruro se sustituye primero.
\end{exercise}

\begin{exercise}[$\star\star$]\label{exo:b3:complex-mechanisms:8}
Enumérense las pruebas experimentales que mostrarían que una transferencia electrónica es de esfera
interna.
\end{exercise}

\begin{exercise}[$\star\star$]\label{exo:b3:complex-mechanisms:9}
Para $\lambda = \qty{80}{kJ/mol}$, calcúlese $\Delta G^{\ddagger}$ para $\Delta G^\circ = 0$, $-20$, $-80$ y \qty{-120}{kJ/mol}.
\end{exercise}

\begin{exercise}[$\star\star\star$]\label{exo:b3:complex-mechanisms:10}
Dos pares tienen las constantes de autointercambio $k_{11} = 4.0$ y $k_{22} = \qty{2.0e5}{L.mol^{-1}.s^{-1}}$, y
su reacción cruzada tiene $K_{12} = \num{1.0e4}$ (datos del ejercicio). Estímese $k_{12}$ ($f = 1$).
\end{exercise}

\begin{exercise}[$\star\star\star$]\label{exo:b3:complex-mechanisms:11}
Con el modelo solo $\sigma$, calcúlense las energías de activación del campo de ligandos de los
iones $d^5$, $d^7$ y $d^8$ de alto espín, y ordénense \ce{Mn^{2+}}, \ce{Co^{2+}} y \ce{Ni^{2+}} según la velocidad
esperada de intercambio de agua.
\end{exercise}

\begin{exercise}[$\star\star\star$]\label{exo:b3:complex-mechanisms:12}
¿Por qué es difícil observar la \omterm{def:b3:complex-mechanisms:inverted}{región invertida} en las reacciones entre iones que
se encuentran por difusión en disolución, y cómo la revelaron las moléculas rígidas dador--aceptor?
\end{exercise}

\section{Problema: fabricar cisplatino, no transplatino}

\begin{problem}[{Problema de fin de semana --- fabricar cisplatino, no transplatino: el
ion platino(II) plano-cuadrado, una síntesis planificada con el efecto trans, la
acuación del fármaco y por qué se activa dentro de las células}]
\label{pb:b3:complex-mechanisms:1}
% ledger: ghs4:K2PtCl4
El cisplatino, $cis$-$\ce{[PtCl2(NH3)2]}$, se fabrica a partir de $\ce{K2[PtCl4]}$. Datos del problema: a
\qty{37}{\celsius}, la primera acuación, $\ce{[PtCl2(NH3)2] + H2O -> [PtCl(H2O)(NH3)2]+ + Cl-}$, tiene la
constante de velocidad $k = \qty{9.0e-5}{s^{-1}}$ y la constante de equilibrio $K = \qty{3.0e-3}{mol/L}$; la
concentración de cloruro es de unos \qty{100}{mM} en el plasma sanguíneo y de \qty{4}{mM} dentro de las células.

\medskip\noindent\textbf{Parte I --- El platino(II).}
\begin{enumerate}
  \item Dese el número de electrones $d$ del \ce{Pt^{2+}}.
  \item ¿Por qué sus complejos son plano-cuadrados y diamagnéticos?
  \item ¿Por qué sus sustituciones son asociativas?
  \item Escríbase la ley de velocidad de dos términos y explíquense sus dos términos.
  \item ¿Es el fármaco \omterm{def:b3:complex-mechanisms:labile}{lábil} o inerte a la escala de tiempo de un día?
\end{enumerate}

\medskip\noindent\textbf{Parte II --- La síntesis.}
\begin{enumerate}[resume]
  \item ¿Qué da el $\ce{K2[PtCl4]}$ con dos \ce{NH3}? ¿Por qué?
  \item Esa vía directa da un producto impuro. La vía de Dhara convierte primero
        el $\ce{[PtCl4]^{2-}}$ en $\ce{[PtI4]^{2-}}$ con KI en exceso. ¿Por qué es útil el yoduro?
  \item La adición de \ce{NH3} da $cis$-$\ce{[PtI2(NH3)2]}$. Explíquese la estereoquímica.
  \item Los yoduros se eliminan después con nitrato de plata en agua. ¿Qué se forma?
  \item La adición de KCl da el cisplatino. ¿Por qué no cambia la configuración?
  \item ¿Cómo se fabricaría, en cambio, el transplatino?
  \item ¿Por qué el transplatino no sirve como fármaco, aunque también reacciona?
\end{enumerate}

\medskip\noindent\textbf{Parte III --- La acuación.}
\begin{enumerate}[resume]
  \item Calcúlese el tiempo de semirreacción de la primera acuación.
  \item Calcúlese la fracción acuada al cabo de \qty{2.0}{h} en una disolución sin cloruro.
  \item ¿Por qué el acuocomplejo es la forma reactiva frente al ADN?
  \item ¿Es la acuación disociativa o asociativa? ¿Cuál sería su \omterm{def:b3:complex-mechanisms:activation-volume}{volumen
        de activación}?
  \item ¿Por qué se prepara en suero salino una disolución del fármaco para inyección?
\end{enumerate}

\medskip\noindent\textbf{Parte IV --- Plasma y célula.}
\begin{enumerate}[resume]
  \item Escríbase la fracción de equilibrio del acuocomplejo en función de
        $[\ce{Cl-}]$.
  \item Calcúlese en el plasma.
  \item Calcúlese dentro de una célula.
  \item Calcúlese la razón.
  \item ¿Por qué, entonces, el fármaco actúa sobre todo dentro de las células?
  \item ¿Qué otros ligandos del interior de una célula pueden unirse al platino y desactivar el
        fármaco?
  \item ¿Por qué los dos ligandos ammina permanecen unidos todo el tiempo?
  \item Enúnciese el resultado: la razón entre la fracción de acuocomplejo dentro de una célula y
        la del plasma.
\end{enumerate}
\end{problem}
